\chapter{Differentialgleichungen}


\section{Fallstudie: mathematisches Pendel}
\begin{figure}[H]
  \begin{center}
    \begin{tikzpicture}[>=latex,scale=1]
      \def\a{25}
      \coordinate (m) at (-90+\a:3);
      \coordinate (maxA) at (-90+1.5*\a:3);
      \coordinate (S) at (0,-3.5);
      \coordinate (Frad) at ($ (0,0)!1.6!(m) $);
      \coordinate (Ftan) at ($ (m)+(\a:3) $);
      \coordinate (Fel) at ($ (m)+(2,0) $);
      \coordinate (Fg) at ($ (m)+(0,-2) $);
      \fill[pattern = north east lines] (-1,0) rectangle (1,0.3);
      \fill (0,0) circle (0.5mm);
      \draw[color=red,dashed] (maxA) arc (-90+1.5*\a:-90-1.5*\a:3);
      \draw[semithick,dotted] (0,0)--(S);
      \shade[ball color=gray!50] (m) circle (1.5mm)
      node[above,shift={(0.2,0.2)}] {$m$};  \draw[thick] (-1,0)--(1,0);
      \markangle{S}{0,0}{m}{79}{\small $\varphi$}
      \draw[thick,blue!51!black] (0,0)--(m) node[midway,right]{$L$};
      \draw[thick,->,black!50!green] (m)--(Fg)
      node[left,shift={(-0,0)}]{$\vec{F}_g=m\vec{g}$};
      \draw[thick,black!50!green,->] (m)--($(m)!(Fg)!(Frad)$)
      node[midway,right]{$\vec{F}_\mathrm{rad}$};
      \draw[semithick,dotted,black!50!green]($(m)!(Fg)!(Ftan)$)--(Fg);
      \draw[semithick,dotted,black!50!green]($(m)!(Fg)!(Frad)$)--(Fg);
      \draw[thick,black!50!green,->] (m)--($(m)!(Fg)!(Ftan)$)
      node[below,inner sep=1]{$\vec{F}_\mathrm{tan}$};
    \end{tikzpicture}
  \end{center}
  \caption{mathematisches Pendel mit Punktmasse mit Masse $m$, Länge $L$ des Seils und Auslenkungswinkel $\varphi$.}\label{fig:pendel}
\end{figure}

Der Betrag der Rückstellkraft steigt mit dem Auslenkungswinkel $\varphi$ bezüglich der Ruhelage. Hierbei zeigt der Vektor der Rückstellkraft $F_{\text{tan}}$ immer in Richtung der Ruheposition, daher ergibt sich ein Minus in folgender Gleichung:
\begin{align*}
  F_{\text{tan}}(t) = -mg\sin(\varphi (t)).
\end{align*}
Beim Betrachten eines schwingenden Fadenpendels zeigt sich, dass die Geschwindigkeit mit zunehmender Auslenkung abnimmt und nach Erreichen des Scheitelpunkts die Richtung wechselt. Die Geschwindigkeitsänderung bedeutet, dass die Pendelmasse eine Beschleunigung erfährt, genauer gesagt findet eine Tangentialbeschleunigung statt, da eine kreisförmige Bewegungsbahn vorliegt. Die Bewegungsgleichung lautet nach dem 2. Newtonschen Gesetz:
\begin{align*}
  m\cdot a_{\text{tan}}(t) = F_{\text{tan}}(t).
\end{align*}

Die Tangentialbeschleunigung lässt sich durch die Winkelbeschleunigung $\ddot{\varphi}$ ausdrücken (zweite zeitliche Ableitung des Winkels nach der Zeit)
\begin{align*}
  a_{\text{tan}}(t) = L\cdot\ddot{\varphi}(t).
\end{align*}
Bei der ungestörten Schwingung stellt die Rückstellkraft des Pendels die einzige äußere Kraft dar. Nach Umstellen erhalten wir:
\begin{align*}
  m\cdot L\cdot \ddot{\varphi}(t) = -m\cdot g\cdot \sin(\varphi (t)).
\end{align*}
Wir stellen fest, dass sich die Masse $m>0$ kürzen lässt. Die Bewegung ist also unabhängig von der Masse des Pendels! Schliesslich erhalten wir:
\begin{align*}
  \ddot{\varphi} (t) + \frac{g}{L}\cdot\sin(\varphi (t)) = 0,
\end{align*}
wobei $L$ die Länge des (masselosen) Seils ist, mit dem das Pendel an der Decke befestigt ist, $g$ ist die Gravitationskonstante und schliesslich ist der Auslenkungswinkel gemäss \cref{fig:pendel}.

